\documentclass[11pt,letterpaper]{article}
\usepackage{fontspec}
\usepackage{tocloft}
\usepackage{multicol}
\usepackage{nicefrac}
\usepackage[left=1.4in,right=1.5in,top=1in,bottom=1in]{geometry}
\setmainfont[Scale=1.4,AutoFakeBold=1.5,AutoFakeSlant=0.3]{Doves Type}

\widowpenalty=10000
\clubpenalty=10000
\interlinepenalty=500

\title{BBQ Sauce}
\author{}
\date{}

\begin{document}

\maketitle
\thispagestyle{empty}

\section*{Ingredients}
\setlength{\columnsep}{20pt}
\begin{multicols}{2}
\noindent
    Braising liquid, defatted \dotfill 1 cup \\
    Ketchup \dotfill \nicefrac{1}{2} cup \\
    Apple cider vinegar \dotfill 2 Tbsp. \\
    Worcestershire sauce \dotfill 1 Tbsp. \\
    Molasses, unsulfured \dotfill 1 Tbsp. \\
    \columnbreak
    Brown sugar, packed \dotfill 1 Tbsp. \\
    Smoked paprika \dotfill 1 tsp. \\
    Dry mustard \dotfill \nicefrac{1}{2} tsp. \\
    Garlic powder \dotfill \nicefrac{1}{2} tsp. \\
    Cayenne pepper \dotfill \nicefrac{1}{4} tsp. \\
    Kosher salt \dotfill \nicefrac{1}{2} tsp. \\
    Black pepper \dotfill \nicefrac{1}{4} tsp. \\
\end{multicols}

\section*{Directions}

\begin{enumerate}
    \item Pour 1~cup defatted \textbf{braising liquid} into a medium saucepan. Bring to a boil over medium-high heat and reduce until approximately \nicefrac{3}{4}~cup remains, about \textit{4--6~minutes}. The liquid will darken and smell intensely of pork and smoke.

    \item Reduce heat to medium. Add \nicefrac{1}{2}~cup \textbf{ketchup}, 2~Tbsp. \textbf{apple cider vinegar}, 1~Tbsp. \textbf{Worcestershire sauce}, 1~Tbsp. \textbf{molasses}, 1~Tbsp. \textbf{brown sugar}, 1~tsp. \textbf{smoked paprika}, \nicefrac{1}{2}~tsp. \textbf{dry mustard}, \nicefrac{1}{2}~tsp. \textbf{garlic powder}, and \nicefrac{1}{4}~tsp. \textbf{cayenne pepper}. Stir to combine.

    \item Simmer over medium-low heat, stirring occasionally, until sauce is glossy and lightly thickened, about \textit{12--15~minutes}. It should coat the back of a spoon but remain pourable.

    \item Taste and adjust: more \textbf{vinegar} for brightness, more \textbf{cayenne} for heat, more \textbf{brown sugar} if the braising liquid was particularly acidic. Add \nicefrac{1}{2}~tsp. \textbf{kosher salt} and \nicefrac{1}{4}~tsp. \textbf{black pepper}, then re-taste and adjust if needed.

    \item Remove from heat. Use immediately for glazing \textbf{ribs}, or cool and transfer to a jar. Makes about 1\nicefrac{1}{2}~cups.
\end{enumerate}

\newpage

{\footnotesize
\setlength{\columnsep}{20pt}
\setlength{\multicolsep}{6pt}
\begin{multicols}{2}
\setlength{\parindent}{0pt}
\setlength{\parskip}{2pt}
\setlength{\itemsep}{0pt}
\setlength{\parsep}{0pt}

\section*{Equipment Required}
\begin{itemize}
    \item Medium saucepan
    \item Silicone spatula or wooden spoon
    \item Glass jar or squeeze bottle (for storage)
    \item Measuring cups and spoons
\end{itemize}

\section*{Hints and Notes}

\subsection*{Yield}
\begin{itemize}
    \item Makes about 1\nicefrac{1}{2}~cups (enough to glaze one full rack of baby back ribs with extra for serving).
\end{itemize}

\subsection*{Mise en Place}
\begin{itemize}
    \item Measure and have ready: 2~Tbsp. \textbf{apple cider vinegar}, 1~Tbsp. \textbf{Worcestershire sauce}, 1~Tbsp. \textbf{molasses}, 1~Tbsp. \textbf{brown sugar}, 1~tsp. \textbf{smoked paprika}, \nicefrac{1}{2}~tsp. \textbf{dry mustard}, \nicefrac{1}{2}~tsp. \textbf{garlic powder}, and \nicefrac{1}{4}~tsp. \textbf{cayenne pepper}
\end{itemize}

\subsection*{Ingredient Tips}
\begin{itemize}
    \item This recipe assumes 1~cup defatted \textbf{braising liquid}; reduce to \nicefrac{3}{4}~cup before adding the remaining ingredients
    \item \textbf{Molasses} (not blackstrap) provides depth and a slight bitterness that anchors the sweetness; blackstrap is too aggressive here
    \item \textbf{Dry mustard} emulsifies the sauce slightly and adds a low, pungent heat distinct from cayenne
    \item \textbf{Ketchup} provides tomato body and a pre-built sweet-acid baseline; use a quality brand (Heinz) rather than off-brand, which tend to be cloyingly sweet
    \item If the braising liquid reduced too far or is very salty, thin the finished sauce with 1--2~Tbsp. water and re-taste before adding more salt
\end{itemize}

\subsection*{Preparation Tips}
\begin{itemize}
    \item Defat the \textbf{braising liquid} thoroughly---pork fat will make the sauce greasy and prevent it from adhering properly to the ribs during air frying
    \item A fat separator is the most efficient tool; lacking one, chill the liquid briefly so the fat solidifies and lifts cleanly
    \item Watch the reduction closely once the liquid drops below 1~cup---it can scorch quickly at this concentration
    \item The sauce will appear thin off the heat; it thickens as it cools---do not over-reduce
    \item For a deeper, more complex sauce, add 2~Tbsp. brewed black coffee with the \textbf{ketchup}; it amplifies the smoky, bitter register without tasting of coffee
\end{itemize}

\subsection*{Make Ahead \& Storage}
\begin{itemize}
    \item Sauce keeps \textit{2~weeks} refrigerated in a sealed glass jar
    \item Reheat gently over low heat or microwave in \textit{30-second} intervals, stirring between each
    \item Sauce thickens considerably when cold; thin with a splash of \textbf{apple cider vinegar} or water when reheating
    \item The sauce base (without \textbf{ketchup} and spices) can be frozen up to \textit{3~months}; finish the sauce fresh after thawing
\end{itemize}

\subsection*{Serving Suggestions}
\begin{itemize}
    \item Primary use: two-coat glaze for the air fryer finish on the \textbf{Pork Baby Back Ribs} (p.~1)
    \item Serve remaining sauce warm on the side in a small bowl or squeeze bottle
    \item Excellent as a dipping sauce for Texas toast or brush on smoked chicken thighs in the last \textit{10~minutes} of cooking
    \item Stir 2~Tbsp. into pulled pork while still hot to sauce the meat from within
    \item Use as a braising sauce for baked beans---add \nicefrac{1}{4}~cup per can of pinto beans with a splash of water and bake at \textit{325°F} for \textit{45~minutes}
\end{itemize}

\end{multicols}
}

\end{document}